\documentclass[11pt,a4paper]{article}

% ── Page geometry ────────────────────────────────────────────────────────────
\usepackage[a4paper, top=2.5cm, bottom=2.5cm, left=2.5cm, right=2.5cm]{geometry}

% ── Fonts & encoding ─────────────────────────────────────────────────────────
\usepackage[T1]{fontenc}
\usepackage[utf8]{inputenc}
\usepackage{lmodern}
\usepackage{microtype}

% ── Hyperlinks ────────────────────────────────────────────────────────────────
\usepackage[colorlinks=true, linkcolor=black, citecolor=black, urlcolor=blue]{hyperref}

% ── Math ─────────────────────────────────────────────────────────────────────
\usepackage{amsmath, amssymb}

% ── Tables ───────────────────────────────────────────────────────────────────
\usepackage{booktabs}
\usepackage{array}

% ── Code listings ────────────────────────────────────────────────────────────
\usepackage{listings}
\usepackage{xcolor}

\definecolor{codebg}{rgb}{0.96,0.96,0.96}
\definecolor{codeframe}{rgb}{0.80,0.80,0.80}

\lstset{
  backgroundcolor=\color{codebg},
  frame=single,
  rulecolor=\color{codeframe},
  basicstyle=\small\ttfamily,
  breaklines=true,
  captionpos=b,
  xleftmargin=1em,
  xrightmargin=1em,
  aboveskip=0.8em,
  belowskip=0.8em,
  showstringspaces=false
}

% ── Section spacing ──────────────────────────────────────────────────────────
\usepackage{titlesec}
\titlespacing*{\section}{0pt}{1.4ex plus 0.5ex}{0.8ex}
\titlespacing*{\subsection}{0pt}{1.0ex plus 0.3ex}{0.5ex}

% ── Paragraph spacing ────────────────────────────────────────────────────────
\setlength{\parskip}{0.4em}
\setlength{\parindent}{0pt}

% ── Abstract width ───────────────────────────────────────────────────────────
\usepackage{abstract}
\renewcommand{\abstractnamefont}{\normalfont\bfseries}

% ─────────────────────────────────────────────────────────────────────────────
\begin{document}

% ── Title block ──────────────────────────────────────────────────────────────
\title{\textbf{QUANTA: Engineering a Production-Ready\\
Post-Quantum Blockchain with Falcon-512 Lattice Signatures}}

\author{
  Kishore K\\
  \small Independent Researcher\\
  \small \href{mailto:kishorekkumar34@gmail.com}{kishorekkumar34@gmail.com}\\
  \small \url{https://www.quantachain.org} \quad
  \url{https://github.com/quantachain/quanta}
}

\date{February 2026}
\maketitle
\thispagestyle{empty}

% ── Abstract ─────────────────────────────────────────────────────────────────
\begin{abstract}
We present QUANTA, an open-source, production-oriented blockchain designed
from genesis to resist attacks from both classical and quantum adversaries.
Unlike hybrid or migration-based approaches, QUANTA integrates two
NIST-standardized post-quantum primitives as its sole cryptographic
foundation: Falcon-512 (FIPS 206) for transaction signatures and
Kyber-1024 (ML-KEM, FIPS 203) for key encapsulation in node-to-node
communication. No classical cryptographic primitive appears in the
critical path. We describe five concrete engineering contributions:
(1)~a strict, domain-separated Falcon-512 verification pipeline that
eliminates known malleability and cross-protocol replay vectors;
(2)~a parallel signature verification engine that reduces full-block
validation time by a factor of six on commodity multi-core hardware;
(3)~an LRU-backed signature cache achieving approximately 80\% hit rates,
nearly eliminating redundant cryptographic work;
(4)~a storage layer combining binary serialization with Zstandard
compression that reduces on-disk block size by 75\% relative to plain JSON
encoding; and
(5)~a crypto-agility mechanism that allows future algorithm migration via a
single-byte scheme tag without a hard fork.
QUANTA is freely available under the MIT licence, with a pre-mainnet testnet
publicly accessible via Docker at
\url{https://hub.docker.com/repository/docker/xd637/quanta-node}.
\end{abstract}

\hrule
\vspace{1em}

% ─────────────────────────────────────────────────────────────────────────────
\section{Introduction}

The resilience of deployed blockchain networks rests almost entirely on the
computational hardness of the elliptic curve discrete logarithm problem
(ECDLP). Bitcoin, Ethereum, and the majority of public chains use ECDSA or
EdDSA; both are asymptotically broken by Shor's algorithm~\cite{shor1994}
running on a sufficiently large fault-tolerant quantum computer. Conservative
engineering analyses project such machines capable of attacking 256-bit
elliptic curve keys within 10--15 years~\cite{mosca2018,aggarwal2018}. The
blockchain setting is particularly acute because transactions and public keys
are permanently on-chain: an attacker harvesting public keys today can
retroactively forge signatures once a quantum computer becomes available.

Existing post-quantum blockchain efforts fall into two broad categories.
The first retrofits PQ signatures onto an existing ECC-based chain via a soft
or hard fork (e.g., Algorand's announced migration research). This leaves
legacy addresses vulnerable during the transition period and requires complex
coexistence logic. The second category consists of purpose-built PQ chains;
the most well-known example is the Quantum Resistant Ledger (QRL), which uses
XMSS~\cite{hulsing2018}, a stateful hash-based scheme that imposes per-key
signing limits and requires careful state management.

QUANTA belongs to the second category but innovates in three ways.
First, it uses Falcon-512~\cite{falcon2020}, a NIST-standardized (FIPS
206~\cite{fips206}) stateless lattice-based scheme, avoiding the key-state
problem of XMSS while providing compact keys and signatures.
Second, it solves the determinism problem that arises when variable-time
Falcon operations interact with multi-architecture consensus: all arithmetic
in the consensus layer is expressed in pure integer math, and the compiler
toolchain is pinned.
Third, it introduces a practical performance architecture---parallel
verification, cached verification, and compressed storage---that makes the
overhead of PQ signatures acceptable in a 10-second-block-time chain.
Critically, QUANTA also employs Kyber-1024 (ML-KEM, FIPS 203) for
post-quantum key encapsulation in all node-to-node communication, ensuring
that both the authentication and confidentiality layers are quantum-resistant
from genesis.

\textbf{Novelty relative to prior work.}
Table~\ref{tab:comparison} summarizes how QUANTA differs from the three closest prior systems.

\begin{table}[h]
\centering
\small
\begin{tabular}{@{}lcccc@{}}
\toprule
\textbf{Property} & \textbf{QRL} & \textbf{QAN} & \textbf{Algorand PQ} & \textbf{QUANTA} \\
\midrule
Signature type & XMSS (stateful) & Dilithium & ECDSA+Falcon hybrid & Falcon-512 (stateless) \\
KEM layer & None & None & None & Kyber-1024 \\
Consensus-level determinism & Not addressed & Not addressed & Not addressed & Explicit (\S\ref{sec:determinism}) \\
Classical fallback & No & Yes & Yes & No \\
Compressed storage & No & No & No & Yes (75\% reduction) \\
Crypto-agility & No & No & No & Yes (tag-dispatch) \\
\bottomrule
\end{tabular}
\caption{Comparison of post-quantum blockchain systems.}
\label{tab:comparison}
\end{table}

QRL's XMSS imposing per-key signing limits and requiring external
state management is a significant operational burden in the blockchain
setting~\cite{hulsing2018}. QAN inherits Ethereum's EVM account model
complexity and does not address cross-architecture consensus determinism.
Algorand's hybrid approach retains classical ECDSA keys in the critical
security path, which is broken by the weakest-link argument: a sufficiently
powerful CRQC can still compromise hybrid-signed transactions by attacking the
ECDSA component. QUANTA eliminates all classical primitives from the critical
path at genesis.

\textbf{Contributions.}
\begin{itemize}
  \item A \textbf{protocol hardening framework} for integrating Falcon-512 into consensus (\S\ref{sec:hardening}).
  \item A \textbf{parallel and cached block-validation engine} achieving 6$\times$ speedup (\S\ref{sec:parallel}).
  \item A \textbf{determinism enforcement model} covering arithmetic, toolchain, and compilation flags (\S\ref{sec:determinism}).
  \item A \textbf{compressed storage architecture} reducing disk costs by 75\% (\S\ref{sec:storage}).
  \item A \textbf{crypto-agility design} enabling future algorithm migration without a hard fork (\S\ref{sec:agility}).
  \item A \textbf{dual-layer PQ security model} pairing Falcon-512 (on-chain) with Kyber-1024 (off-chain), providing the first complete PQ stack with neither layer relying on a classical fallback (\S\ref{subsec:kyberimpl}).
\end{itemize}

% ─────────────────────────────────────────────────────────────────────────────
\section{Background}

\subsection{Falcon-512}

Falcon~\cite{falcon2020} is a lattice-based digital signature scheme based on
NTRU lattices and Fast Fourier sampling. It was selected by NIST in 2022 and
standardized as FIPS~206. At security level~1 (Falcon-512), parameters are:

\begin{table}[h]
\centering
\begin{tabular}{@{}ll@{}}
\toprule
\textbf{Parameter} & \textbf{Size} \\
\midrule
Public key  & 897 bytes \\
Secret key  & $\approx$1281 bytes \\
Signature   & 33--666 bytes (variable, compressed) \\
\bottomrule
\end{tabular}
\caption{Falcon-512 key and signature sizes.}
\end{table}

The variable-length compressed signature is a practical challenge for
blockchain engineering: block size bounds cannot be computed from transaction
counts alone. We address this in \S\ref{subsec:blocksize}.

Falcon has no known quantum attacks; the best known classical attack requires
$2^{128}$ operations (comparable to AES-128). The scheme is stateless, making
it directly applicable to the blockchain setting without additional state
management.

\subsection{Threat Model}
\label{subsec:threat}

We consider an adversary with access to a cryptographically relevant quantum
computer (CRQC) capable of running Shor's algorithm. The adversary targets:
(i)~signature forgery on past transactions whose public keys are on-chain;
(ii)~consensus manipulation via forged block headers; and (iii)~cross-chain
and cross-protocol signature replay. We do not consider attacks on
proof-of-work (QUANTA uses SHA3-256 PoW, which is Grover-resistant with a
128-bit classical security margin after the square-root speedup).

\subsection{Kyber-1024 (ML-KEM)}
\label{subsec:kyber}

Kyber~\cite{fips203} is a lattice-based key encapsulation mechanism (KEM)
based on the Module Learning With Errors (MLWE) problem. It was standardized
by NIST as FIPS~203 (ML-KEM) in 2024. At security level~5 (Kyber-1024),
it provides 256-bit post-quantum security with the following parameters:

\begin{table}[h]
\centering
\begin{tabular}{@{}ll@{}}
\toprule
\textbf{Parameter} & \textbf{Size} \\
\midrule
Public key  & 1568 bytes \\
Secret key  & 3168 bytes \\
Ciphertext  & 1568 bytes \\
Shared secret & 32 bytes \\
\bottomrule
\end{tabular}
\caption{Kyber-1024 (ML-KEM) key and ciphertext sizes.}
\end{table}

In QUANTA, Kyber-1024 is used for encrypted peer-to-peer node communication.
When two nodes establish a connection, they perform a Kyber key encapsulation
handshake to derive a shared session key, which then encrypts all traffic
between them. This renders the full network stack quantum-resistant: Falcon-512
secures the authenticity of on-chain data (transactions and blocks), while
Kyber-1024 secures the confidentiality of off-chain node communication.
Neither primitive has a classical fallback.

% ─────────────────────────────────────────────────────────────────────────────
\section{Protocol Hardening for Falcon-512}
\label{sec:hardening}

The most subtle challenge in deploying Falcon-512 in a consensus system is the
gap between what the reference library guarantees (correct cryptographic
verification) and what consensus requires (deterministic, malleability-free,
replay-resistant verification). We address this through four mechanisms.

\subsection{Domain-Separated Canonical Signing}

Every transaction signature in QUANTA covers not the raw serialized
transaction bytes but the hash
$\texttt{SHA3-256}(\textit{``QUANTA\_TX\_V1:''} \;\|\; \textit{signing\_bytes})$.
The prefix \textit{``QUANTA\_TX\_V1:''} is a domain separator baked into the
protocol at genesis; it is a consensus-critical constant that must never
change without a hard fork.

\begin{lstlisting}[caption={Canonical signing hash construction.}]
// Domain prefix ensures no cross-protocol signature reuse.
SIGNING_DOMAIN = "QUANTA_TX_V1:"

canonical_signing_hash(data):
    hasher = SHA3_256.new()
    hasher.update(SIGNING_DOMAIN)
    hasher.update(data)
    return hasher.finalize()   // 32-byte output
\end{lstlisting}

This serves two purposes. First, it prevents cross-protocol signature reuse: a
valid QUANTA signature cannot be replayed in any other protocol because the
hash input is different. Second, fixing the Falcon message to exactly 32 bytes
eliminates message-length variance in the signed-message encoding, removing a
class of length-extension attacks. The test suite includes an explicit case
that verifies a signature produced with the domain prefix does not verify
against the raw SHA3 hash of the same data, confirming domain separation is
enforceable at runtime.

\subsection{Pre-Verification Size Guards}

Falcon's reference binding returns an error for malformed inputs, but parsing
malformed data still consumes CPU time and, under network load, can become a
denial-of-service bottleneck. QUANTA rejects inputs before entering Falcon
internals using strict byte-length pre-checks:

\begin{lstlisting}[caption={Strict length bounds enforced before cryptographic verification.}]
FALCON512_PUBKEY_BYTES  = 897   // exact; consensus-critical constant
FALCON512_SIG_MAX_BYTES = 698   // 666-byte sig + 32-byte message
FALCON512_SIG_MIN_BYTES = 33    // minimum plausible

verify_strict(message, signed_msg, public_key):
    if len(public_key) != FALCON512_PUBKEY_BYTES:  return False
    if not (SIG_MIN <= len(signed_msg) <= SIG_MAX): return False
    // cryptographic verification follows ...
\end{lstlisting}

These bounds are documented as consensus-critical. A node that accepts
signatures outside these bounds would diverge from the network. The maximum
signed-message size of 698 bytes is derived from Falcon's maximum signature
length (666 bytes) plus the fixed 32-byte message digest.

\subsection{Sender Binding}

Address derivation in QUANTA computes the first 20 bytes of
$\texttt{SHA3-256}(\textit{public\_key})$, prefixed with \texttt{0x} in
hexadecimal representation. The verification function enforces that the public
key embedded in a transaction hashes to the declared sender address
\emph{before} cryptographic verification runs. This prevents key-substitution
attacks in which an adversary replaces the public key with one for which they
hold the secret key while keeping the sender address unchanged---a class of
attack that pure Falcon verification would not detect.

\subsection{Memory Safety for Secret Keys}

The keypair abstraction annotates secret key storage with automatic
zeroing-on-drop semantics. Secret key material is overwritten with zeros when
the keypair object leaves scope, preventing recovery via heap inspection, core
dumps, or swap files. This is a standard hardening practice, but is worth
documenting explicitly in a blockchain context where consensus nodes may share
a memory address space with other processes.

\subsection{Security Analysis}
\label{sec:security}

We argue informally that the four mechanisms in \S\ref{sec:hardening} jointly
prevent the attack classes described in the threat model (\S\ref{subsec:threat}).

\noindent\textbf{Cross-protocol replay (Claim~1).}
An adversary who holds a valid Falcon-512 signature $\sigma$ over message $m$
in a distinct protocol $\Pi'$ cannot replay $\sigma$ in QUANTA. The signed
message in QUANTA is $\texttt{SHA3-256}(\texttt{``QUANTA\_TX\_V1:''} \;\|\; m)$;
the signed message in $\Pi'$ is some distinct hash (lacking the domain
prefix). These values are distinct with probability $1 - 2^{-256}$ by
collision resistance of SHA3-256. QUANTA's verification function
rejects any signature that does not verify against the domain-prefixed hash,
so replay from $\Pi'$ fails at the Falcon verification step.

\noindent\textbf{Key-substitution attacks (Claim~2).}
Suppose an adversary knows a Falcon-512 keypair $(sk', pk')$ and targets
a victim address $\mathit{addr}_v = \texttt{SHA3-256}(pk_v)[0..20]$.
If $pk' \neq pk_v$, then
$\texttt{SHA3-256}(pk')[0..20] \neq \texttt{SHA3-256}(pk_v)[0..20]$
with overwhelming probability (collision resistance). The sender-binding check
rejects the transaction before Falcon verification is called.

\noindent\textbf{Denial-of-service via malformed signatures (Claim~3).}
Parsing and verifying a malformed Falcon signed-message requires $O(|\sigma|)$
work in the reference implementation. Pre-verification length guards reject
any input with $|pk| \neq 897$ or $|\text{signed\_msg}| \notin [33, 698]$
in $O(1)$, reducing the attacker's forcing cost for DoS to the cost of
producing correctly-sized but otherwise invalid inputs, which Falcon's
internal checks then reject in bounded time.

\noindent\textbf{Secret key exfiltration via memory inspection (Claim~4).}
The keypair object's destructor overwrites secret key storage with zeros
before releasing memory. An adversary who can read heap or swap pages after
the keypair leaves scope observes only zeros, not the Falcon secret key.
This provides defence-in-depth; it does not substitute for process isolation
or OS-level security.

\noindent\textbf{Scheme downgrade (Claim~5).}
Each transaction carries a scheme tag byte. The dispatcher hard-rejects
unrecognised tag values (\S\ref{sec:agility}). An adversary cannot cause
a QUANTA node to evaluate a signature under a weaker scheme by manipulating
the tag: any unknown tag is rejected outright, and the classical-scheme tag
space is left permanently unassigned.

% ─────────────────────────────────────────────────────────────────────────────
\section{Parallel and Cached Block Validation}
\label{sec:parallel}

Falcon-512 signature verification takes approximately 1.5\,ms per operation on
a modern x86-64 core. A block containing 1,200 transactions (the QUANTA
maximum) would require 1,800\,ms for serial verification---well above the
10-second block target. We address this with two complementary optimizations.

\subsection{Parallel Signature Verification}

Block validation distributes signature verification across all available CPU
cores using a data-parallel work-stealing scheduler. Each worker thread
independently verifies one transaction's Falcon-512 signature with no shared
mutable state:

\begin{lstlisting}[caption={Data-parallel block validation; correctness follows from pure-function semantics of Falcon verify.}]
// All transactions verified in parallel; no shared mutable state.
all_valid = block.transactions
    .parallel_iter()
    .all(|tx| tx.is_system_tx() OR tx.verify_signature())
\end{lstlisting}

On an 8-core machine this reduces signature verification time from
$\approx$1,800\,ms to $\approx$225\,ms---a 6$\times$ speedup. The correctness
argument is straightforward: signature verification is a pure function, so
data-parallel evaluation is safe. Balance and nonce validation, which requires
sequential state, follows in a separate serial pass.

\subsection{LRU Signature Verification Cache}

Transactions propagate across the P2P network before appearing in a block. A
consensus node that verified a transaction's signature during mempool admission
would otherwise verify it again at block arrival. A bounded Least-Recently-Used
(LRU) cache with 100,000 entries avoids this redundant work. The cache maps
each transaction's 256-bit hash to a boolean verification result; on a hit the
Falcon computation is skipped entirely. In practice, with a mempool that
overlaps significantly with block content, the hit rate exceeds 80\%.

\begin{table}[h]
\centering
\begin{tabular}{@{}lc@{}}
\toprule
\textbf{Strategy} & \textbf{1,200-tx block validation time} \\
\midrule
Serial, no cache                    & $\approx$1,800\,ms \\
Parallel (8 cores), no cache        & $\approx$225\,ms   \\
Parallel + 80\% cache hit rate      & $\approx$45\,ms    \\
\bottomrule
\end{tabular}
\caption{Combined effect of parallelism and caching on block validation latency.}
\end{table}

\subsection{Block Size Accounting for Variable Signatures}
\label{subsec:blocksize}

Because Falcon-512 signatures are variable-length (33--666 bytes), block
construction cannot rely on a fixed per-transaction size. QUANTA's block
template builder serializes each candidate transaction and measures the
resulting byte count before adding it to the candidate set, maintaining a
running accumulator. Transactions that would cause the accumulator to exceed
the 2\,MB block size limit are skipped even if the transaction count limit has
not yet been reached.

% ─────────────────────────────────────────────────────────────────────────────
\section{Build and Consensus Determinism}
\label{sec:determinism}

A signature scheme's correctness is necessary but not sufficient for consensus
safety. If two honest nodes produce different verification results for the same
transaction because their binaries were compiled differently, the network will
fork. We enforce determinism at three levels.

\subsection{Toolchain Pinning}

The Rust compiler version is pinned to 1.85.0 for all nodes via a
repository-level toolchain specification file. Nodes are explicitly prohibited
from using nightly or beta compiler channels. A nightly LLVM back-end may
apply different loop vectorization passes to the Falcon C library, potentially
changing execution timing or, in edge cases, floating-point rounding in
auxiliary computations, leading to consensus divergence.

\subsection{Integer-Only Consensus Arithmetic}

All consensus-critical computations---block reward calculation, difficulty
adjustment---use pure 64-bit unsigned integer arithmetic. The annual reward
reduction illustrates the principle: rather than computing
$\text{reward} \times 0.85$, the implementation multiplies by 85 and divides
by 100 iteratively:

\[
  \text{reward} \;\leftarrow\; \lfloor \text{reward} \times 85 \;/\; 100 \rfloor
\]

Using IEEE~754 double-precision floating-point here would be a consensus bug.
The x87 FPU on x86-64 produces intermediate results at 80-bit extended
precision, whereas ARM64 strict-IEEE mode produces 64-bit results; the two can
differ in the last bit of the mantissa. A disagreement in computed reward or
difficulty causes one node to reject another's blocks, producing a consensus
fork across architectures. The integer approach accumulates less than 0.01
QUA error over 20 years, well within the 5 QUA minimum reward floor.

Difficulty adjustment uses the same strategy: the ratio
$\text{expected\_time} / \text{actual\_time}$ is scaled by 1,000, rounded by
adding 500, then divided---entirely in unsigned 64-bit arithmetic with explicit
overflow guards.

\subsection{Reproducible Release Profile}

The release build profile enforces three settings for byte-for-byte binary
reproducibility: (i)~a single code-generation unit, which eliminates LLVM's
non-deterministic inlining decisions that arise when multiple compilation units
are merged in parallel; (ii)~link-time optimization (LTO), which performs
cross-module optimization deterministically at link time; and (iii)~disabled
incremental compilation, because incremental caches accumulate machine-specific
metadata that must not contaminate consensus-node binaries.

% ─────────────────────────────────────────────────────────────────────────────
\section{Compressed Storage Architecture}
\label{sec:storage}

Falcon-512 transactions carry 897-byte public keys and up to 666-byte
signatures in every record. A na\"{i}ve JSON-serialized block at 1,200
transactions occupies approximately 3--4\,MB on disk. To make full-node
operation feasible on commodity hardware, QUANTA implements a two-stage
compression pipeline.

\subsection{Binary Serialization}

JSON uses variable-length ASCII encoding with per-field structural overhead
(names, braces, quotes). Replacing JSON with a compact binary serialization
format reduces the serialized size by approximately 22\% and improves
throughput by approximately $8\times$ (no UTF-8 encoding, no escape handling).

\subsection{Zstandard Block Compression}

After binary serialization, each block is compressed with Zstandard
(level~3) before being written to the embedded key-value store.
Level~3 offers an excellent speed-to-ratio trade-off: on blockchain
data---highly repetitive address strings and signature byte patterns---Zstandard
achieves compression ratios of 3--4$\times$. Combined with binary
serialization, the overall reduction relative to plain JSON is approximately
75\%. A full-archive node therefore consumes roughly 28\% of the disk space of
a na\"{i}ve implementation.

\subsection{On-Demand Block Loading}

The in-memory blockchain representation keeps only the genesis block resident.
All other blocks are loaded from the persistent store on demand. This reduces
steady-state RAM usage from $O(h)$ (chain height) to $O(1)$:

\begin{itemize}
  \item \textbf{Before:} 3.15M blocks $\times$ 2\,MB $\approx$ 6.3\,TB RAM (infeasible)
  \item \textbf{After:} genesis block only $\approx$ 1\,MB RAM
\end{itemize}

A 1,000-block LRU cache sits in front of the disk store, making common
operations---get latest block, validate recent transactions---efficient without
materializing the full chain.

\subsection{Transaction Indexing}

Without secondary indexing, finding a transaction by hash requires a sequential
scan of all blocks---$O(h)$ in chain height. At year-one height this would take
approximately 30\,seconds. QUANTA writes a secondary index at block-save time
that maps each transaction hash to a \emph{(block height, position within block)}
tuple. Lookup is therefore $O(1)$: resolve the tuple from the index, load the
enclosing block (likely cached), and return the transaction at the recorded
position. This reduces lookup time from $\approx$30\,s to $\approx$10\,ms.


% ─────────────────────────────────────────────────────────────────────────────
\section{Crypto Agility}
\label{sec:agility}

Despite Falcon-512's strong security guarantees, prudent design mandates a
clear upgrade path in case of future cryptanalytic advances. QUANTA embeds a
single-byte algorithm tag in every transaction. The current tag space defines:

\begin{table}[h]
\centering
\begin{tabular}{@{}cll@{}}
\toprule
\textbf{Tag} & \textbf{Algorithm} & \textbf{Status} \\
\midrule
\texttt{0x01} & Falcon-512  & Active \\
\texttt{0x02} & ML-DSA (Dilithium) & Reserved \\
\texttt{0x03} & SLH-DSA (SPHINCS+) & Reserved \\
\bottomrule
\end{tabular}
\caption{Algorithm tag space in the QUANTA wire format.}
\end{table}

During verification, the consensus layer dispatches on the algorithm tag
before calling any cryptographic function. An unknown tag value causes the
transaction to be hard-rejected, preventing silent acceptance of future-scheme
transactions by unupgraded nodes. A new algorithm can be activated by
assigning it the next available tag and passing a soft fork: upgraded nodes
accept both old and new schemes; unupgraded nodes reject the new scheme but
continue to accept the old one. No wire-format changes are required.

The domain separator includes an explicit version token, allowing future
scheme versions to use a distinct prefix without collision, so signatures under
different schemes cannot be replayed across scheme versions. This design was
informed by NIST's recommendation to plan for algorithm migration when
deploying PQ schemes in long-lived systems~\cite{fips206}.

\subsection{M-of-N Post-Quantum Multi-Signatures}

QUANTA implements M-of-N threshold signing over Falcon-512. A multi-signature
transaction record stores $N$ Falcon-512 public keys, the threshold $M$, and a
slot vector of $N$ optional signature entries. Each collected signature is
verified against the canonical signing hash before being stored, preventing
accumulation of invalid signatures. The multi-signature address is derived as
the first 20 bytes of
$\texttt{SHA3-256}(pk_1 \;\|\; pk_2 \;\|\; \cdots \;\|\; pk_N)$,
cryptographically binding the authorized key set to the address.

% ─────────────────────────────────────────────────────────────────────────────
\section{Implementation and Deployment}
\label{sec:implementation}

\subsection{Language Choice: Rust}

QUANTA is implemented entirely in Rust ($\approx$9,000 lines of source, 2021
edition). The choice of Rust over C, C++, or Go is deliberate and motivated by
three properties critical for a security-sensitive consensus system:

\begin{enumerate}
  \item \textbf{Memory safety without a garbage collector.} Rust's ownership
        and borrow-checker model eliminates entire classes of memory-safety
        bugs---buffer overflows, use-after-free, dangling pointers---at compile
        time, without the unpredictable pause latency of a tracing GC. For a
        10-second block-time chain, GC pauses would introduce jitter into
        consensus timing.
  \item \textbf{Data-race freedom.} Rust's type system statically prevents data
        races across threads. The parallel signature verification engine
        (\S\ref{sec:parallel}) shares no mutable state between worker threads;
        Rust enforces this at compile time.
  \item \textbf{Predictable performance.} Rust compiles to native machine code
        with zero-cost abstractions, producing throughput and latency
        characteristics comparable to C---essential for keeping Falcon-512
        verification within the block budget without hardware acceleration.
\end{enumerate}

\subsection{Post-Quantum Network Layer (Kyber-1024)}
\label{subsec:kyberimpl}

While Falcon-512 secures on-chain data, all off-chain communication between
QUANTA nodes is encrypted using Kyber-1024 (ML-KEM). When two peers connect,
they execute a Kyber key encapsulation handshake to derive a shared 32-byte
session secret, which bootstraps a symmetric cipher for the connection
lifetime. This design has two important properties:

\begin{itemize}
  \item \textbf{No classical fallback.} The handshake does not offer
        ECDH or any classical key agreement variant. A connecting peer that
        does not support Kyber-1024 is rejected, preventing downgrade attacks.
  \item \textbf{Forward secrecy.} A fresh Kyber keypair is generated for
        each session. Compromise of a long-term key does not expose past
        session traffic.
\end{itemize}

The result is a blockchain where both layers of the security stack---the
on-chain authenticity layer (Falcon-512) and the off-chain confidentiality
layer (Kyber-1024)---are fully quantum-resistant, with no cryptographic
primitive that a CRQC can break.

\subsection{Runtime Stack}

The runtime uses an asynchronous I/O scheduler for networking, an HTTP
framework for the REST API (port 3000), and a custom TCP-based P2P protocol
(port 8333). A JSON-RPC daemon (port 7782) enables operator control.

\begin{table}[h]
\centering
\begin{tabular}{@{}ll@{}}
\toprule
\textbf{Component} & \textbf{Library / Standard} \\
\midrule
Async runtime       & Tokio 1.35 \\
PQ signatures       & pqcrypto-falcon 0.3.0 (version-pinned) \\
PQ key encapsulation & pqcrypto-kyber 0.8 (Kyber-1024) \\
Hashing             & SHA3-256 (FIPS 202) \\
Storage             & sled 0.34 (embedded key-value store) \\
Compression         & Zstandard 0.13 \\
Parallel compute    & Work-stealing thread pool \\
Verification cache  & LRU cache, 100k entries \\
Serialization       & Compact binary encoding \\
\bottomrule
\end{tabular}
\caption{QUANTA runtime stack.}
\end{table}

The Falcon-512 library dependency is pinned to an exact version. Any version
change alters the compiled Falcon C code and could change signature encoding,
requiring a coordinated network-wide upgrade.

\subsection{Wallet Security}

HD wallet derivation uses BIP39 24-word mnemonics with Argon2id key
stretching. Wallet files encrypt the Falcon-512 secret key with
ChaCha20-Poly1305 using an Argon2id-derived key, providing memory-hard
brute-force resistance against offline attacks.


\subsection{Public Testnet}

A pre-mainnet testnet node is publicly accessible via Docker:

\begin{lstlisting}
docker pull xd637/quanta-node:latest
docker run -d -p 3000:3000 -p 8333:8333 xd637/quanta-node
\end{lstlisting}

The full source code is available under the MIT licence at
\url{https://github.com/quantachain/quanta}.

% ─────────────────────────────────────────────────────────────────────────────

\section{Discussion and Limitations}

\textbf{Proof-of-Work.} QUANTA uses SHA3-256 Adaptive Proof-of-Work. PoW is
simple to reason about and provides Sybil resistance without additional trust
assumptions, but carries well-known environmental concerns. Future protocol
versions may explore proof-of-stake variants with post-quantum verifiable
random functions.

\textbf{Signature size overhead.} At 666 bytes maximum, Falcon-512 signatures
are substantially larger than ECDSA (71 bytes) or Ed25519 (64 bytes). Our
compression layer (\S\ref{sec:storage}) partially offsets this. For
applications where bandwidth dominates, Dilithium is a trade-off worth
evaluating, and the crypto-agility design (\S\ref{sec:agility}) makes such a
migration feasible.

\textbf{Falcon timing variation.} The Falcon signing algorithm uses randomized
Gaussian sampling whose running time is variable. In a blockchain context,
signing happens only in wallets, not in consensus nodes, so timing
side-channels in signing do not affect verification determinism. Verification
is a deterministic function.

\textbf{Testnet status.} The current deployment is a pre-mainnet testnet.
Production readiness requires external security audits, extended stress
testing, and peer review of the consensus parameter choices.

% ─────────────────────────────────────────────────────────────────────────────
\section{Related Work}

\textbf{QRL}~\cite{hulsing2018} implements XMSS, a stateful hash-based scheme
standardized in RFC~8391. XMSS provides strong security with smaller signatures
than Falcon but requires per-key signing state, complicating wallet management
and key recovery.

\textbf{QAN Platform} deploys Kyber and Dilithium on a layer-1 compatible with
Ethereum's EVM. This preserves EVM tool compatibility but inherits EVM's
account model complexity; QAN does not address consensus-level determinism
challenges in the depth that QUANTA does.

\textbf{Algorand's PQ research} focuses on hybrid signatures (ECDSA + Falcon)
for migration, accepting a transient period where classical keys coexist with
quantum-resistant keys. QUANTA takes the position that hybrid security is
weaker than pure PQ security by the weakest-link argument.

\textbf{Ethereum's roadmap}~\cite{buterin2023} acknowledges eventual PQ
migration but has not committed to a concrete algorithm or timeline, leaving
existing account keys permanently at risk once CRQCs arrive.

% ─────────────────────────────────────────────────────────────────────────────
\section{Conclusion}

We presented QUANTA, an open-source blockchain built exclusively on
NIST-standardized post-quantum cryptography. Our contributions address the gap
between academic PQ scheme correctness and the engineering requirements of a
production consensus system: strict domain-separated verification prevents
cross-protocol replay; parallel verification and LRU caching make 1,200-transaction
block validation feasible within a 10-second block budget; integer-only
consensus arithmetic eliminates floating-point non-determinism across
architectures; and Zstandard compression reduces on-disk block size by 75\%.
The crypto-agility mechanism provides a clear, non-disruptive migration path
to future algorithms as the PQ landscape matures.

QUANTA demonstrates that post-quantum blockchains are not merely academically
feasible but practically engineerable today, with performance characteristics
competitive with classical-crypto chains on commodity hardware.

% ─────────────────────────────────────────────────────────────────────────────
\begin{thebibliography}{10}

\bibitem{shor1994}
P.~W. Shor,
``Algorithms for quantum computation: discrete logarithms and factoring,''
\textit{Proc.\ 35th Annual Symposium on Foundations of Computer Science (FOCS)},
pp.~124--134, 1994.

\bibitem{mosca2018}
M.~Mosca,
``Cybersecurity in an era with quantum computers: will we be ready?''
\textit{IEEE Security \& Privacy}, vol.~16, no.~5, pp.~38--41, 2018.

\bibitem{aggarwal2018}
D.~Aggarwal, G.~K. Brennen, T.~Lee, M.~Santha, and M.~Tomamichel,
``Quantum attacks on Bitcoin, and how to protect against them,''
\textit{Ledger}, vol.~3, 2018.

\bibitem{hulsing2018}
A.~Hülsing, D.~Butin, S.-L. Gazdag, J.~Rijneveld, and A.~Mohaisen,
``XMSS: Extended Merkle Signature Scheme,''
RFC 8391, Internet Engineering Task Force, 2018.

\bibitem{falcon2020}
P.-A. Fouque et al.,
``Falcon: Fast-Fourier Lattice-based Compact Signatures over NTRU,''
NIST PQC Round~3 Submission, 2020.
Available: \url{https://falcon-sign.info/}

\bibitem{fips203}
National Institute of Standards and Technology,
``Module-Lattice-Based Key-Encapsulation Mechanism Standard,''
FIPS~203, 2024.

\bibitem{fips206}
National Institute of Standards and Technology,
``Module-Lattice-Based Digital Signature Standard,''
FIPS~206, 2024.

\bibitem{buterin2023}
V.~Buterin,
``The Ethereum roadmap: post-quantum migration,''
Ethereum Research Forum, 2023.
Available: \url{https://ethresear.ch/}

\bibitem{quanta2026}
K.~K.,
``QUANTA: Engineering a Production-Ready Post-Quantum Blockchain with Falcon-512 Lattice Signatures,''
Zenodo, Feb.\ 2026.
DOI: \href{https://doi.org/10.5281/zenodo.18753528}{10.5281/zenodo.18753528}.
Available: \url{https://github.com/quantachain/quanta}

\bibitem{quantadocker}
K.~K.,
``QUANTA Node Docker Image,''
Docker Hub, 2026.
Available: \url{https://hub.docker.com/repository/docker/xd637/quanta-node}

\bibitem{pornin2022}
T.~Pornin,
``New Efficient, Constant-Time Implementations of Falcon,''
\textit{IACR Transactions on Cryptographic Hardware and Embedded Systems}, 2022.

\end{thebibliography}

\end{document}
